\textbf{Время:} 02:00

\textbf{Тезис:} столкновение ядер позволяет исследовать расширение плотной системы многих взаимодействующих частиц.

\textbf{Визуал:} начальное столкновение → плотное облако → расширение; связать рисунок с вопросами о размере и движении.

\textbf{Текст:} Тяжёлое ядро содержит много нуклонов — протонов и нейтронов. Сталкивая такие ядра, мы создаём много взаимодействующих частиц в малом объёме. Часть энергии движения превращается в возбуждение и рождение новых частиц. Система расширяется, становится разреженной и распадается на конечные частицы.

Мы хотим понять, насколько она успела увеличиться, как двигалось вещество и как частицы его покидали. Разные модели эволюции могут предсказывать разные размеры области испускания, даже если похожим образом описывают число конечных частиц. Значит, размеры дают дополнительную проверку физической картины.

В нашем рассказе будем использовать пионы — лёгкие адроны, составные частицы сильного взаимодействия. У пиона спин ноль; положительный и отрицательный пионы — разные сорта. Их много, а пары одинакового заряда удобны для исследования корреляций. При высокой энергии возможно образование кварк-глюонной плазмы; для применимости самого метода это не обязательное исходное условие.

Пауза: вернуться к двум вопросам — насколько большая система и как она расширялась?

\textbf{Переход:} Прежде чем измерять размер, уточним, о какой стадии столкновения идёт речь.
